\section{Uniform completion of the near-vertex regime}
\label{sec:vertex-completion}

The fixed-count hypothesis in Theorem~\ref{thm:simplex-fixed-vertex}
can be removed at the level of the leading crossing threshold. The
result is uniform over every rare composition with total mass $o(b)$.
An accompanying example explains why the corresponding probability
approximation need not be uniform: its error can diverge even while
the relative error of the crossing threshold tends to zero.

Throughout this section, $\ell\ge1$ is fixed, $a_1,\ldots,a_\ell$
are positive integers, $A=\sum_i a_i$, and
\[
 H(y)=\prod_{i=1}^{\ell}B_{a_i}(y),\qquad
 \kappa(y)=\frac{yH'(y)}{H(y)}.
\]
Write $y_a$ for the unique positive solution of $H(y_a)=1$.

\begin{lemma}[elasticity and concavity of the rare-count polynomial]
\label{lem:vertex-elasticity}
The polynomial $H$ is convex, $\log H$ is concave on $(0,\infty)$,
and $\kappa$ is nondecreasing. Moreover,
\begin{equation}\label{eq:vertex-elasticity}
 \frac1{A+1}\le y_a\le1,\qquad
 \kappa(y)\le \ell+\sqrt{\ell Ay},
 \qquad \kappa(y)\ge \max\{\ell,\tfrac12\sqrt{Ay}\}
 \quad(0<y\le1).
\end{equation}
\end{lemma}

\begin{proof}
Convexity follows directly from the nonnegative coefficients. For one
factor, give $J\in\{0,\ldots,a-1\}$ probabilities proportional to
\[
 r_j=\binom{a-1}{j}\frac{y^j}{(j+1)!}.
\]
The successive ratio
$R(j)=(j+1)r_{j+1}/r_j=y(a-1-j)/(j+2)$, with $R(a-1)=0$,
is nonincreasing. Summing the ratio identity gives
$\mathbb E J=\mathbb E R(J)$ and
$\mathbb E J^2=\mathbb E[(J+1)R(J)]$. Therefore
\[
 \operatorname{Var}J=\mathbb E J+\operatorname{Cov}(J,R(J))
 \le \mathbb E J.
\]
The covariance is nonpositive because, for an independent copy $J'$,
twice the covariance equals
$\mathbb E[(J-J')(R(J)-R(J'))]\le0$.
Putting $M=J+1$ and $\kappa_a=yB_a'/B_a=\mathbb E M$, we obtain
\[
 (\log B_a)''=\frac{\operatorname{Var}M-\mathbb E M}{y^2}\le0,
 \qquad y\kappa_a'=\operatorname{Var}M\ge0.
\]
This proves both qualitative assertions without invoking the locations
of the zeros of a Laguerre polynomial.

The elementary differential equation $yB_a''+yB_a'=aB_a$ gives
\[
 ay=\kappa_a^2+\operatorname{Var}M+(y-1)\kappa_a.
\]
Using $0\le\operatorname{Var}M\le\kappa_a-1$ yields
$\kappa_a\le1+\sqrt{ay}$ and
$\kappa_a^2+y\kappa_a\ge ay+1$.
For $0<y\le1$, these imply
$\kappa_a\ge\max\{1,\sqrt{ay}-1/2\}
\ge\tfrac12\sqrt{ay}$.
Summing over the factors and applying Cauchy--Schwarz proves the
elasticity bounds. Finally, $B_a(1)\ge1$ and
\[
 B_a(y)\le y\Phi(ay)\le y e^{ay/2},
\]
where $(j+1)!\ge2^j$ proves the last inequality term by term.
Thus $H(1)\ge1$, whereas
$H(1/(A+1))\le(A+1)^{-\ell}e^{A/(2(A+1))}<1$.
\end{proof}

\begin{lemma}[an exact compound-binomial representation]
\label{lem:vertex-randomization}
For $0<u<1$, put $v=u/(1-u)$. Let $G$ have the gamma distribution
with shape $2$ and rate $1$, let $E_1,\ldots,E_{b-1}$ be independent
rate-one exponential variables, and let $Z_1,\ldots,Z_{b-1}$ be
independent Bernoulli variables with parameter $u$. Take all these
variables independent and set
\[
 T=G+\sum_{j=1}^{b-1}Z_jE_j.
\]
Then
\begin{equation}\label{eq:vertex-randomization}
 S_{(b,a)}(u)=(1-u)^A\mathbb E[H(vT)],\qquad
 \mathbb E e^{sT}=(1-s)^{-2}
 \left(1+\frac{us}{1-s}\right)^{b-1}\quad(s<1).
\end{equation}
In particular, set $y=bu^2$, $\kappa=\kappa(y)$ and
$s=\kappa/[bu(1-u)]$. Whenever $s<1$,
\begin{align}
 (1-u)^A H(y)&\le S_{(b,a)}(u),\label{eq:vertex-jensen-lower}\\
 S_{(b,a)}(u)&\le (1-u)^A H(y)e^{-\kappa}(1-s)^{-2}
 \left(1+\frac{us}{1-s}\right)^{b-1}.
 \label{eq:vertex-tangent-upper}
\end{align}
\end{lemma}

\begin{proof}
The exact positive-refresh integral is
\[
 S_{(b,a)}(u)=\frac{(1-u)^{b+A-1}}v
 \int_0^\infty e^{-t}B_b(vt)H(vt)\,dt.
\]
The integral of $e^{-t}B_b(vt)$ is $v(1+v)^{b-1}$.
After normalization its density is a mixture of gamma densities with
shape $J+2$, where $J\sim\operatorname{Bin}(b-1,u)$. This is exactly
the distribution of $T$, proving the identity and its moment generating
function. Since $\mathbb ET=2+(b-1)u$,
$v\mathbb ET\ge bu^2=y$. Convexity and monotonicity of $H$ prove
\eqref{eq:vertex-jensen-lower}. Concavity of $\log H$ gives, for $w>0$,
\[
 H(w)\le H(y)\exp\!\left(\frac{\kappa}{y}(w-y)\right).
\]
Insert $w=vT$ and use the moment generating function to prove
\eqref{eq:vertex-tangent-upper}.
\end{proof}

\begin{theorem}[all growing rare counts near a vertex]
\label{thm:vertex-uniform-root}
For fixed $\ell$, uniformly over all positive rare counts with $A=o(b)$,
the unique positive crossing odds satisfy
\begin{equation}\label{eq:vertex-uniform-root}
 u_{b,a}=\sqrt{\frac{y_a}{b}}
 \left(1+O_\ell\!\left(\sqrt{\frac{A+1}{b}}\right)\right),
 \qquad \prod_i B_{a_i}(y_a)=1.
\end{equation}
A quantitative, unoptimized version is as follows. Put
\[
 K=\ell+\sqrt\ell+1,\quad
 \varepsilon=\sqrt{\frac{A+1}{b}},\quad \delta=16K\varepsilon.
\]
If $\delta\le1/2$, then
\begin{equation}\label{eq:vertex-uniform-bracket}
 (1-\delta)\sqrt{y_a/b}\ \le u_{b,a}\le
 (1+\delta)\sqrt{y_a/b}.
\end{equation}
No restriction such as $A(\log A)^2=o(b)$ is needed.
\end{theorem}

\begin{proof}
The proper-run polynomial has positive coefficients, zero constant
coefficient, and is unbounded at infinity; its crossing is unique.
Put $u_0=\sqrt{y_a/b}$ and $\kappa_0=\kappa(y_a)$.
At $u_+=(1+\delta)u_0$ we have $u_+<1/4$.
Monotonicity of $\kappa$ and the lower bound in
Lemma~\ref{lem:vertex-elasticity} give
\[
 \log H(bu_+^2)\ge2\kappa_0\log(1+\delta),\qquad
 \frac{Au_0}{\kappa_0}\le2\sqrt{A/b}\le2\varepsilon.
\]
Consequently, by~\eqref{eq:vertex-jensen-lower},
\[
 \log S(u_+)\ge
 2\kappa_0\log(1+\delta)-\frac{Au_+}{1-u_+}
 \ge\kappa_0\left(\frac43\delta-4\varepsilon\right)>0.
\]

At $u_-=(1-\delta)u_0$, write $y=bu_-^2$ and
$\kappa=\kappa(y)$. The inequalities $y_a\ge1/(A+1)$,
$\delta\le1/2$, and~\eqref{eq:vertex-elasticity} imply
\[
 u_-\le\varepsilon,\qquad bu_-\ge\frac1{2\varepsilon},\qquad
 s=\frac{\kappa}{bu_-(1-u_-)}
 \le 2(2\ell+\sqrt\ell)\varepsilon\le4K\varepsilon\le\frac18.
\]
Using $\log(1+x)\le x$, $-\log(1-s)\le s/(1-s)$ and
$u_-,s\le1/4$ in~\eqref{eq:vertex-tangent-upper}, and dropping the
nonpositive term $A\log(1-u_-)$, gives
\begin{align*}
 \log\frac{S(u_-)}{H(y)}
 &\le -\kappa+\frac{2s}{1-s}+\frac{bu_-s}{1-s}\\
 &=\kappa\left(\frac1{(1-u_-)(1-s)}-1\right)
       +\frac{2s}{1-s}\\
 &\le 5\kappa(u_-+s)\le25K\kappa\varepsilon.
\end{align*}
Here $\kappa\ge\ell\ge1$ was used in the last two inequalities.
Again by monotonicity of $\kappa$,
\[
 \log H(y)\le2\kappa\log(1-\delta)\le-2\kappa\delta.
\]
Thus $\log S(u_-)\le\kappa(-32K+25K)\varepsilon<0$.
The two strict inequalities and uniqueness bracket the crossing.
\end{proof}

\begin{corollary}[an explicit shape law for growing rare counts]
\label{cor:vertex-shape-law}
Suppose $A\to\infty$, $A=o(b)$, and the actual proportions
$\alpha_i=a_i/A$ stay in a fixed compact subset of the positive
$\ell$-simplex. Put $h=\sum_i\sqrt{\alpha_i}$. Then
\[
 u_{b,a}\sim\frac{\ell\log A}{2h\sqrt{Ab}}.
\]
More precisely, the polynomial root satisfies
\begin{align*}
 2h\sqrt{Ay_a}
 ={}&\ell\log A-\frac\ell2\log\log A
 +\frac\ell2\log\!\left(\frac{8\pi h}{\ell}\right)
 +\frac34\sum_i\log\alpha_i
 +O_\ell\!\left(\frac{\log\log A}{\log A}
                  +\frac{(\log A)^2}{A}\right).
\end{align*}
Replacing $\sqrt{Ay_a}$ on the left by $\sqrt{Ab}\,u_{b,a}$ incurs
an additional $O_\ell(\log A\sqrt{A/b})$ error.
\end{corollary}

\begin{proof}
The Laguerre--Bessel inequality, uniformly when $y\to0$, gives
$H(y)=y^\ell\prod_i\Phi(a_i y)(1+O_\ell(y))$.
Write $x=\sqrt{Ay}$. Bracketing first at fixed small and large
multiples of $\log A$ shows that the root has $x\asymp\log A$.
The standard large-argument expansion of $I_1$ then gives uniformly
over the stated proportions
\[
 H(x^2/A)=
 \frac{x^{\ell/2}e^{2hx}}
 {A^\ell(4\pi)^{\ell/2}(\prod_i\alpha_i)^{3/4}}
 \left(1+O_\ell(x^{-1}+x^2/A)\right).
\]
Taking logarithms, substituting
$x=(\ell/(2h))\log A+O_\ell(\log\log A)$ in the logarithmic term,
and using Theorem~\ref{thm:vertex-uniform-root} proves the claims.
The constants also depend on the fixed compact set of proportions.
\end{proof}

\begin{proposition}[probability approximation can fail while roots agree]
\label{prop:vertex-probability-failure}
Let $a=\lfloor b/\log b\rfloor$, let $y_a$ solve $B_a(y_a)=1$,
and set $u_0=\sqrt{y_a/b}$. Then
\[
 H(y_a)=B_a(y_a)^2=1,\qquad
 S_{(b,a,a)}(u_0)\longrightarrow\infty,
 \qquad \frac{u_{b,a,a}}{u_0}\longrightarrow1.
\]
Thus crossing roots admit a uniform approximation in a regime where the
corresponding relative probability approximation fails:
$S(u)\sim H(bu^2)$ is false as a uniform assertion over $A=o(b)$.
\end{proposition}

\begin{proof}
Retain only run words that start and end in the large color and have
exactly one adjacent pair of distinct rare colors. The run enumeration
in Theorem~\ref{thm:simplex-fixed-vertex} shows that their contribution
is the polynomial $2uyB_a'(y)^2$, with each term of total rare-run
count $m$ additionally multiplied by
$\prod_{j=1}^{m-1}(1-j/b)_+$.
Give $J$ probabilities proportional to
$\binom{a-1}{j}y^j/j!$, the coefficients of $B_a'(y)$.
Its differential equation
$y(B_a')''+(1+y)(B_a')'-(a-1)B_a'=0$ implies
$\mathbb E J^2+y\mathbb E J=(a-1)y$; hence
$\mathbb E J\le\sqrt{ay}$ and $\mathbb E J^2\le ay$.
For two independent copies, $m=J_1+J_2+2$, so
\[
 \mathbb E[m(m-1)]\le4ay+6\sqrt{ay}+2.
\]
The clipped-product deficit bound therefore proves
\[
 S_{(b,a,a)}(u_0)\ge
 \frac{2\kappa_a(y_a)^2}{bu_0}
 \left(1-\frac{4ay_a+6\sqrt{ay_a}+2}{2b}\right).
\]
The preceding corollary for one rare color gives
$\sqrt{ay_a}\sim\tfrac12\log a$, and the elasticity bounds give
$\kappa_a(y_a)\sim\sqrt{ay_a}$.
The last display is asymptotic to a lower bound
$\sqrt{a/b}\log a\sim\sqrt{\log b}$, which diverges.
The root assertion follows directly from
Theorem~\ref{thm:vertex-uniform-root}, since $2a=o(b)$.
\end{proof}

These results complete the leading near-vertex threshold for any fixed
number of rare colors. They do not classify every finite-size global
mode, nor do they assert a uniform higher-order root expansion for all
growing rare counts. The analytic arguments in this section have
written proofs and are not included in the finite Lean formalization.
